% Generated draft: requires mathematical and source review.
\documentclass[11pt]{article}
\usepackage[margin=25mm]{geometry}
\usepackage{fontspec}
\setmainfont{Latin Modern Roman}
\usepackage{amsmath,amssymb,mathtools,unicode-math}
\setmathfont{Latin Modern Math}
\usepackage{enumitem,needspace,xurl,hyperref}
\setlength{\parindent}{0pt}
\setlength{\parskip}{5pt}
\setlength{\emergencystretch}{3em}
\newcommand{\N}{\mathbb{N}}
\newcommand{\Z}{\mathbb{Z}}
\newcommand{\Q}{\mathbb{Q}}
\newcommand{\R}{\mathbb{R}}
\newcommand{\C}{\mathbb{C}}
\newcommand{\F}{\mathbb{F}}
\newcommand{\A}{\mathbb{A}}
\newcommand{\PP}{\mathbb{P}}
\newcommand{\E}{\mathbb{E}}
\newcommand{\eps}{\varepsilon}
\newcommand{\dd}{\,\mathrm{d}}
\begin{document}
\section*{UnsolvedMath ID: 800239}
\textbf{Source problem code:} catalogue entry
\subsection{Definitions and mathematical statement}
Let \(n\geq 1\), and let
\[
B^n=\left\{z\in\C^n:\sum_{j=1}^n|z_j|^2<1\right\}.
\]
Let \(\Omega\subset\C^n\) be a bounded domain for which there exists a biholomorphic map \(F:B^n\to\Omega\), meaning that \(F\) is bijective and both \(F\) and \(F^{-1}\) are holomorphic. No regularity of \(\partial\Omega\), or extension of \(F\) to the boundary, is assumed. The operator \(\bar\partial\) is the antiholomorphic exterior derivative; on scalar functions it is
\[
\bar\partial u=\sum_{j=1}^n\frac{\partial u}{\partial\bar z_j}\,\mathrm d\bar z_j.
\]
The source asks whether \(\bar\partial\) satisfies uniform estimates on every such \(\Omega\), emphasizing domains with very rough boundary.

The phrase ``uniform estimates'' is not defined in the supplied record. Consequently, the source does not determine a unique fully quantified estimate. For example, if the intended meaning is supremum norm solvability for bounded \((0,1)\)-forms, one possible precise question is whether there exists \(C_\Omega<\infty\) such that every form
\[
f=\sum_{j=1}^n f_j\,\mathrm d\bar z_j,
\qquad f_j\in L^\infty(\Omega),\qquad \bar\partial f=0
\]
admits \(u\in L^\infty(\Omega)\) satisfying
\[
\bar\partial u=f,\qquad
\|u\|_{L^\infty(\Omega)}\leq C_\Omega\max_{1\leq j\leq n}\|f_j\|_{L^\infty(\Omega)}.
\]
Here the differential equations would be interpreted in the sense of distributions. This displayed formulation is an explicitly conditional interpretation: the source does not specify these data spaces, this form degree, this notion of solution, or this dependence of the constant. It also does not require a bounded linear solution operator. These choices require mathematical and source review before adopting a definitive statement.
\subsection{Short English statement}
Does every bounded domain biholomorphic to a complex unit ball admit uniform estimates for solving \(\bar\partial\)? The source highlights very rough boundaries, but leaves the norms, form degrees, solution spaces, and uniformity of the constants unspecified. The displayed supremum norm formulation is only a possible interpretation.
\subsection{Sources}
\begin{itemize}
\item Local UnsolvedMath registry: permanent ID \(800239\), exact source problem code \texttt{catalogue entry}, title ``Several Complex Variables Problem List — Problem 36.'' The supplied registry labels the problem open; this label is not an independently verified mathematical status assessment.
\item Buzzard, Fornaess, and Krantz, \emph{The Several Complex Variables Problem List} (1995), Problem 36. Supplied source URL: \url{https://arxiv.org/abs/math/9509202}. The registry names these authors in its proposer field; the extracted problem itself is attributed to B. Berndtsson, with contact \texttt{bob@math.chalmers.se}.
\item Supplied provenance records access on 2026-10-16 and extraction by \texttt{arxiv\_latex}. These are supplied source-processing details; no independent source visit or literature verification is asserted here. Supplied upstream revision: \texttt{d371ee917ac9647fbb06cd6fc529deec6f3dda82}.
\end{itemize}
\end{document}
